\uuid{yzcp}
\titre{Oscillateur harmonique et solution fondamentale causale}

\niveau{M1}
\module{Analyse}
\chapitre{Distributions}
\sousChapitre{Solutions fondamentales et convolution}
\theme{Équation de l'oscillateur harmonique}
\auteur{Quentin Liard}
\datecreate{ 2026-09-03}
\organisation{CentraleSupélec}
\difficulte{3}

\contenu{

\texte{
On considère l'équation de l'oscillateur harmonique
$$
T''+\omega^2T=B,
$$
où $B$ représente le signal extérieur et $\omega\in\mathbb{R}$. On cherche d'abord une solution $G\in\mathcal{D}'_+$ de
$$
G''+\omega^2G=\delta.
$$
}

\begin{enumerate}

\item \question{Résoudre le problème de Cauchy
$$
\begin{cases}
z''+\omega^2z=0,\\
z(0)=0,\\
z'(0)=1.
\end{cases}
$$}

\reponse{
On distingue deux cas.

Si $\omega\neq0$, les solutions de l'équation homogène s'écrivent
$$
z(t)=A\cos(\omega t)+B\sin(\omega t).
$$
La condition $z(0)=0$ donne $A=0$, puis $z'(0)=1$ donne $B=1/\omega$. Ainsi,
$$
\boxed{z(t)=\frac{\sin(\omega t)}{\omega}}.
$$

Si $\omega=0$, l'équation devient $z''=0$. Les conditions initiales donnent
$$
\boxed{z(t)=t}.
$$
}

\item \question{En déduire une solution de
$$
G''+\omega^2G=\delta,
$$
puis une solution de
$$
T''+\omega^2T=B.
$$}

\reponse{
Pour une fonction $z$ de classe $C^2$ telle que $z(0)=0$ et $z'(0)=1$, on a
$$
(Hz)''=Hz''+z'(0)\delta+z(0)\delta'=Hz''+\delta.
$$
Par conséquent,
$$
(Hz)''+\omega^2Hz
=H(z''+\omega^2z)+\delta
=\delta.
$$
Ainsi, une solution causale fondamentale est
$$
G(t)=
\begin{cases}
\displaystyle H(t)\frac{\sin(\omega t)}{\omega}, & \omega\neq0,\\[1ex]
H(t)t, & \omega=0.
\end{cases}
$$
Sous l'hypothèse que la convolution avec $B$ est définie, par exemple si $B\in\mathcal{D}'_+$, une solution de l'équation forcée est
$$
\boxed{T=G\ast B}.
$$
En effet,
$$
T''+\omega^2T
=(G''+\omega^2G)\ast B
=\delta\ast B
=B.
$$
}

\end{enumerate}

}
